
\chapter{Metodologia: Ensemble de Clusterização e Regras de Associação}
\label{cap:metodologia}

O pipeline (\texttt{ml/clusterizacao\_regras\_associacao.py}) roda em duas
etapas sobre o Data Mart (\autoref{sec:datamart}): primeiro clusteriza os
municípios, depois descobre regras de associação \textbf{dentro} de cada
cluster, separadamente.

\section{Etapa 1: Clusterização dos municípios}
\label{sec:clusterizacao}

Cada município com volume mínimo de ocorrências ganha um \textbf{perfil de
risco}, agregado de \texttt{fato\_acidentes} + \texttt{dim\_local} +
\texttt{dim\_classificacao\_acidente}, com nove \textit{features}: volume
(em log), taxa de letalidade, taxa de feridos graves, veículos médios,
proporções de fim de semana/noite/chuva/pista simples, e concentração da
causa principal. As \textit{features} são padronizadas
(\texttt{StandardScaler}) antes do agrupamento, para que nenhuma domine a
distância euclidiana só por estar numa escala maior.

O agrupamento usa \textbf{KMeans} \cite{macqueen1967} (\texttt{scikit-learn}
\cite{pedregosa2011}), com $k$ escolhido automaticamente: o pipeline testa
$k$ de 2 a 8 e mantém o valor com o maior \textit{silhouette score}
\cite{rousseeuw1987}, que compara, para cada município, a distância média
ao próprio cluster com a distância média ao cluster vizinho mais próximo,
sem depender de rótulo verdadeiro. O \autoref{code:clusterizacao} reproduz o
núcleo dessa escolha, de \texttt{escolher\_k\_e\_clusterizar()}; o modelo
final é reajustado com o $k$ escolhido.

\begin{center}
  \begin{codigo}[H]
    \small
    \lstinputlisting[language=Python, firstline=216, lastline=227]{../app/ml/clusterizacao_regras_associacao.py}
    \caption{Escolha de $k$ por silhueta}
    \label{code:clusterizacao}
  \end{codigo}
\end{center}

\section{Etapa 2: Regras de associação por cluster}
\label{sec:regras}

As ocorrências de cada cluster viram \textit{cestas}: até 9 itens por
ocorrência, um por classificação (\texttt{causa=...},
\texttt{tipo\_acidente=...}, \texttt{clima=...}, \texttt{tipo\_pista=...}
etc.); valores não informativos são descartados da cesta. Sobre as cestas de
cada cluster roda o \textbf{Apriori} \cite{agrawal1994} (\texttt{mlxtend}
\cite{raschka2018}), que encontra os itemsets frequentes e as regras
antecedente~$\rightarrow$~consequente, filtradas e ordenadas por três
métricas: \textbf{suporte} (frequência conjunta), \textbf{confiança}
($P(\text{consequente} \mid \text{antecedente})$) e \textbf{\textit{lift}}
(razão entre confiança observada e esperada sob independência;
$\textit{lift}>1$ indica associação positiva).

O \autoref{code:regras} reproduz o núcleo dessa etapa: cada cluster vira uma
matriz \textit{one-hot} de cestas, alimenta o Apriori e, por fim, as regras.

\begin{center}
  \begin{codigo}[H]
    \small
    \lstinputlisting[language=Python, linerange={307-309,318-324}]{../app/ml/clusterizacao_regras_associacao.py}
    \caption{Apriori e regras de associação, por cluster}
    \label{code:regras}
  \end{codigo}
\end{center}

O \autoref{code:clusterizacao} e o \autoref{code:regras} são excertos; o
código completo está em \texttt{app/ml/clusterizacao\_regras\_associacao.py}.
Como o \textit{lift} é simétrico, o pipeline mantém só a direção de maior
confiança de cada par.

Parâmetros desta execução: ano de 2024 (um ano por vez, pois os 20 anos
juntos excedem a memória do Apriori por cluster), mínimo de 30 ocorrências
por município, $k$ entre 2 e 8, suporte mínimo 0,05, \textit{lift} mínimo
1,1, confiança mínima 0,30, itemsets de no máximo 2 itens (regras de 1
antecedente para 1 consequente, evitando \textit{supersets} redundantes do
mesmo par) e as 15 regras de maior \textit{lift} mantidas por cluster.

\section{Persistência e disponibilização para o BI}
\label{sec:persistencia}

Os resultados são gravados no próprio Data Mart, nas tabelas
\texttt{cluster\_municipio} e \texttt{regra\_associacao}
(\autoref{fig:dw-estrela}), e em CSV/PNG em \texttt{ml/output/}. O pipeline
roda pela \sigla{DAG}{Directed Acyclic Graph}
\texttt{ml\_clusterizacao\_regras\_associacao} do Airflow, logo após a carga
inicial do DW.
